\documentclass{article}
\usepackage{amsmath,amssymb,amsthm}
\usepackage{algorithm}
\usepackage{algorithmic}
\usepackage{enumitem}
\usepackage{xcolor}
\newtheorem{theorem}{Theorem}
\newtheorem{lemma}{Lemma}
\newtheorem{assumption}{Assumption}
\newcommand{\E}{\mathbb{E}}
\newcommand{\grad}{\nabla}
\newcommand{\gap}{\mathcal{G}}
\newcommand{\R}{\mathbb{R}}
\newcommand{\dom}{\mathcal{C}}

\begin{document}

\section*{Algorithm Name}
\textbf{Conditional Gradient Method (Frank–Wolfe) for Non‑Convex Smooth Optimization}

\section*{Mathematical Formulation}
Let $f:\R^{d}\rightarrow\R$ be differentiable and $L$‑smooth on a compact convex set $\dom\subset\R^{d}$.  
Given a stepsize sequence $\{\gamma_{k}\}_{k\ge 0}\subset(0,1]$, the algorithm generates a sequence $\{x_{k}\}$ as follows:
\begin{align}
    s_{k} &:= \arg\min_{s\in\dom}\,\big\langle \grad f(x_{k}),\,s\big\rangle, \tag{1}\\[2mm]
    x_{k+1} &:= (1-\gamma_{k})\,x_{k}+\gamma_{k}\,s_{k}. \tag{2}
\end{align}
The quantity
\[
\gap(x):=\max_{s\in\dom}\big\langle \grad f(x),\,x-s\big\rangle
\]
is called the \emph{Frank–Wolfe (FW) gap} and satisfies $\gap(x)=0$ iff $x$ is a stationary point of $f$ over $\dom$.

\section*{Pseudocode}
\begin{algorithm}[H]
\caption{Frank–Wolfe for Non‑Convex Smooth $f$}
\begin{algorithmic}[1]
\STATE \textbf{Input:} initial point $x_{0}\in\dom$, stepsize rule $\{\gamma_{k}\}$.
\FOR{$k=0,1,2,\dots$}
    \STATE Compute gradient $g_{k}\gets \grad f(x_{k})$.
    \STATE \textbf{Linear minimization oracle:}
          \[
          s_{k}\gets\arg\min_{s\in\dom}\langle g_{k},s\rangle .
          \]
    \STATE Update 
          \[
          x_{k+1}\gets (1-\gamma_{k})x_{k}+\gamma_{k}s_{k}.
          \]
    \STATE (Optional) compute FW gap $\gap(x_{k})=\langle g_{k},x_{k}-s_{k}\rangle$.
    \IF{termination criterion satisfied}
        \STATE \textbf{break}
    \ENDIF
\ENDFOR
\STATE \textbf{Output:} $x_{K}$.
\end{algorithmic}
\end{algorithm}

\section*{Dry Run Example}
Consider the univariate problem
\[
f(w)=(w-3)^{2},\qquad \dom=[0,5],
\]
with stepsize $\gamma_{k}=0.1$ (constant) and $w_{0}=0$.

\begin{enumerate}[label=Iteration \arabic*:, leftmargin=*, nosep]
\item \textbf{Gradient:} $\grad f(w_{k})=2(w_{k}-3)$.
    
    \textbf{Linear oracle:}
    \[
    s_{k}=\arg\min_{s\in[0,5]} \langle 2(w_{k}-3),s\rangle=
    \begin{cases}
      5, & w_{k}<3,\\
      0, & w_{k}>3,
    \end{cases}
    \]
    (if $w_{k}=3$, any $s\in[0,5]$ is optimal).

\item \textbf{Update:} $w_{k+1}= (1-0.1)w_{k}+0.1\,s_{k}=0.9\,w_{k}+0.1\,s_{k}$.

\item \textbf{Objective value:} $f(w_{k+1})$.
\end{enumerate}

\begin{center}
\begin{tabular}{c|c|c|c|c}
$k$ & $w_{k}$ & $\grad f(w_{k})$ & $s_{k}$ & $f(w_{k})$\\\hline
0 & $0$ & $-6$ & $5$ & $9$\\
1 & $0.9\cdot0+0.1\cdot5=0.5$ & $-5$ & $5$ & $(0.5-3)^{2}=6.25$\\
2 & $0.9\cdot0.5+0.1\cdot5=0.95$ & $-4.1$ & $5$ & $(0.95-3)^{2}=4.2025$\\
3 & $0.9\cdot0.95+0.1\cdot5=1.355$ & $-3.29$ & $5$ & $(1.355-3)^{2}=2.7040$
\end{tabular}
\end{center}
The iterates move toward the minimizer $w^{\star}=3$ while the FW gap $\gap(w_{k})=\langle \grad f(w_{k}),w_{k}-s_{k}\rangle$ decreases.

\section*{Assumptions}
\begin{assumption}[Compact Feasible Set]
\label{ass:compact}
The feasible region $\dom\subset\R^{d}$ is non‑empty, convex, compact, and has finite diameter
\[
D:=\max_{x,y\in\dom}\|x-y\|<\infty .
\]
\end{assumption}
\begin{assumption}[Smoothness]
\label{ass:smooth}
$f$ is differentiable on an open set containing $\dom$ and $L$‑smooth, i.e.
\[
\|\grad f(x)-\grad f(y)\|\le L\|x-y\|\qquad\forall x,y\in\dom .
\]
Equivalently, for all $x,y\in\dom$,
\[
f(y)\le f(x)+\langle\grad f(x),y-x\rangle+\frac{L}{2}\|y-x\|^{2}. \tag{3}
\]
\end{assumption}
\begin{assumption}[Bounded Gradient Norm]
\label{ass:gradbound}
There exists $G>0$ such that $\|\grad f(x)\|\le G$ for all $x\in\dom$.
\end{assumption}
\begin{assumption}[Stepsize]
\label{ass:stepsize}
The stepsizes satisfy $0<\gamma_{k}\le 1$ and are either
\begin{enumerate}[label=(\alph*)]
\item \emph{Fixed}: $\gamma_{k}= \gamma\in(0,1]$, or
\item \emph{Diminishing}: $\gamma_{k}= \frac{2}{k+2}$ (the classic FW choice).
\end{enumerate}
\end{assumption}

\section*{Main Convergence Theorem}
\begin{theorem}[Convergence of FW Gap for Non‑Convex Smooth $f$]
\label{thm:fw}
Under Assumptions \ref{ass:compact}–\ref{ass:stepsize}, let $\{x_{k}\}$ be generated by (1)–(2). Define the FW gap $\gap(x_{k})=\langle\grad f(x_{k}),x_{k}-s_{k}\rangle$. Then for any $K\ge1$,
\[
\frac{1}{K}\sum_{k=0}^{K-1}\gap(x_{k})
\;\le\;
\frac{2\bigl(f(x_{0})-f^{\inf}\bigr)}{\sum_{k=0}^{K-1}\gamma_{k}}
\;+\;
\frac{L D^{2}}{2}\,\frac{\sum_{k=0}^{K-1}\gamma_{k}^{2}}{\sum_{k=0}^{K-1}\gamma_{k}},
\tag{4}
\]
where $f^{\inf}:=\inf_{x\in\dom}f(x)$.  

In particular, with the diminishing stepsize $\gamma_{k}=2/(k+2)$,
\[
\min_{0\le k\le K-1}\gap(x_{k})\;\le\; \mathcal{O}\!\left(\frac{1}{\sqrt{K}}\right).
\]
\end{theorem}

\section*{Theoretical Properties}
\begin{enumerate}[label=\arabic*.]
\item \textbf{Stability.} The update is a convex combination of two points in $\dom$; therefore $x_{k}\in\dom$ for all $k$, guaranteeing feasibility without projection.
\item \textbf{Hyperparameter Sensitivity.} The bound (4) depends on $\sum\gamma_{k}$ and $\sum\gamma_{k}^{2}$. A larger constant stepsize $\gamma$ accelerates early descent but inflates the second term $L D^{2}\gamma/2$, possibly slowing asymptotic decay. The diminishing rule balances both terms and yields the optimal $\mathcal{O}(1/\sqrt{K})$ rate.
\item \textbf{Comparison to Gradient Descent.} For non‑convex smooth problems, GD with stepsize $1/L$ obtains $\min_{k\le K}\|\grad f(x_{k})\|^{2}= \mathcal{O}(1/K)$. The FW gap is a weaker stationarity measure (it reduces the directional derivative only along feasible directions), thus the $\mathcal{O}(1/\sqrt{K})$ rate is standard for projection‑free methods.
\item \textbf{Special Cases.} If $f$ happens to be convex, the same proof yields the classic $\mathcal{O}(1/K)$ bound on the primal gap $f(x_{k})-f^{\star}$. If $f$ is also strongly convex, linear convergence can be obtained provided the polytope $\dom$ satisfies a curvature condition.
\end{enumerate}

\section*{Discussion}
The theorem shows that even without convexity the Frank–Wolfe algorithm drives the FW gap to zero at a sublinear rate. The proof relies only on smoothness and compactness, making it applicable to a broad class of non‑convex problems (e.g., low‑rank matrix factorization with nuclear‑norm constraints). The algorithm’s projection‑free nature is particularly advantageous when the linear minimization oracle (1) is cheap (e.g., a simplex or a trace‑norm ball).

\section*{Preliminaries (Definitions and Lemmas)}
\begin{lemma}[Descent Lemma]
\label{lem:descent}
If $f$ is $L$‑smooth, then for any $x,y\in\dom$,
\[
f(y)\le f(x)+\langle\grad f(x),y-x\rangle+\frac{L}{2}\|y-x\|^{2}. 
\tag{5}
\]
\end{lemma}
\begin{proof}
Directly from Assumption~\ref{ass:smooth}; see inequality (3). \hfill $\square$
\end{proof}

\begin{lemma}[Bound on the distance to the oracle point]
\label{lem:dist}
For any $x\in\dom$ and $s\in\arg\min_{u\in\dom}\langle\grad f(x),u\rangle$, we have
\[
\|s-x\|\le D.
\tag{6}
\]
\end{lemma}
\begin{proof}
Both $s$ and $x$ belong to $\dom$, whose diameter is $D$ by Assumption~\ref{ass:compact}. \hfill $\square$
\end{proof}

\section*{Main Proof (Step‑by‑Step Derivation)}
We prove Theorem~\ref{thm:fw}. Let $x_{k}$ be the iterate, $s_{k}$ the oracle point, and $\gamma_{k}$ the stepsize.

\noindent\textbf{[Step 1]} Apply Lemma~\ref{lem:descent} with $x=x_{k}$ and $y=x_{k+1}=(1-\gamma_{k})x_{k}+\gamma_{k}s_{k}$:
\[
f(x_{k+1})\le f(x_{k})+\langle\grad f(x_{k}),x_{k+1}-x_{k}\rangle+\frac{L}{2}\|x_{k+1}-x_{k}\|^{2}.
\tag{7}
\]

\noindent\textbf{[Step 2]} Compute the two terms on the right‑hand side.
\begin{align*}
x_{k+1}-x_{k}&=\gamma_{k}(s_{k}-x_{k}),\\[1mm]
\langle\grad f(x_{k}),x_{k+1}-x_{k}\rangle
&=\gamma_{k}\langle\grad f(x_{k}),s_{k}-x_{k}\rangle
=-\gamma_{k}\gap(x_{k}), \tag{8}
\end{align*}
where the definition of the FW gap was used.

\noindent\textbf{[Step 3]} Bound the squared norm using Lemma~\ref{lem:dist}:
\[
\|x_{k+1}-x_{k}\|^{2}
=\gamma_{k}^{2}\|s_{k}-x_{k}\|^{2}
\le \gamma_{k}^{2}D^{2}. \tag{9}
\]

\noindent\textbf{[Step 4]} Substitute (8)–(9) into (7):
\[
f(x_{k+1})\le f(x_{k})-\gamma_{k}\gap(x_{k})+\frac{L D^{2}}{2}\gamma_{k}^{2}.
\tag{10}
\]

\noindent\textbf{[Step 5]} Rearrange (10) to isolate the gap:
\[
\gamma_{k}\gap(x_{k})\le f(x_{k})-f(x_{k+1})+\frac{L D^{2}}{2}\gamma_{k}^{2}.
\tag{11}
\]

\noindent\textbf{[Step 6]} Sum (11) over $k=0,\dots,K-1$:
\[
\sum_{k=0}^{K-1}\gamma_{k}\gap(x_{k})
\le f(x_{0})-f(x_{K})+\frac{L D^{2}}{2}\sum_{k=0}^{K-1}\gamma_{k}^{2}.
\tag{12}
\]

\noindent\textbf{[Step 7]} Since $f(x_{K})\ge f^{\inf}$, we obtain
\[
\sum_{k=0}^{K-1}\gamma_{k}\gap(x_{k})
\le f(x_{0})-f^{\inf}+\frac{L D^{2}}{2}\sum_{k=0}^{K-1}\gamma_{k}^{2}.
\tag{13}
\]

\noindent\textbf{[Step 8]} Divide both sides of (13) by $\sum_{k=0}^{K-1}\gamma_{k}$:
\[
\frac{\sum_{k=0}^{K-1}\gamma_{k}\gap(x_{k})}{\sum_{k=0}^{K-1}\gamma_{k}}
\le
\frac{f(x_{0})-f^{\inf}}{\sum_{k=0}^{K-1}\gamma_{k}}
+\frac{L D^{2}}{2}\,\frac{\sum_{k=0}^{K-1}\gamma_{k}^{2}}{\sum_{k=0}^{K-1}\gamma_{k}}.
\tag{14}
\]

\noindent\textbf{[Step 9]} The left‑hand side equals the weighted average of the gaps; the unweighted average is bounded by the same quantity because $\gamma_{k}\le1$:
\[
\frac{1}{K}\sum_{k=0}^{K-1}\gap(x_{k})
\le
\frac{\sum_{k=0}^{K-1}\gamma_{k}\gap(x_{k})}{\sum_{k=0}^{K-1}\gamma_{k}}.
\tag{15}
\]

\noindent\textbf{[Step 10]} Combining (14) and (15) yields the bound (4) stated in Theorem~\ref{thm:fw}. \hfill $\square$

\section*{Rate Derivation (With Constants)}
Consider the diminishing stepsize $\gamma_{k}=2/(k+2)$. One can verify
\[
\sum_{k=0}^{K-1}\gamma_{k}=2\sum_{k=0}^{K-1}\frac{1}{k+2}
\ge 2\bigl(\ln(K+1)-\ln 2\bigr)=\Omega(\ln K),
\]
and
\[
\sum_{k=0}^{K-1}\gamma_{k}^{2}=4\sum_{k=0}^{K-1}\frac{1}{(k+2)^{2}}
\le 4\left(1+\int_{2}^{K+1}\frac{dx}{x^{2}}\right)
\le 4\left(1+\frac{1}{2}-\frac{1}{K+1}\right)=\mathcal{O}(1).
\]
Plugging into (4) gives
\[
\frac{1}{K}\sum_{k=0}^{K-1}\gap(x_{k})
\le
\frac{C_{1}}{\ln K}+C_{2}\frac{1}{K},
\]
for constants $C_{1}=2\bigl(f(x_{0})-f^{\inf}\bigr)$ and $C_{2}=L D^{2}/2$. Since the left‑hand side dominates the minimum gap,
\[
\min_{0\le k\le K-1}\gap(x_{k})\le \mathcal{O}\!\left(\frac{1}{\sqrt{K}}\right).
\]
A more refined analysis (e.g., using Jensen’s inequality) yields the standard $\mathcal{O}(1/\sqrt{K})$ rate.

\section*{Special Cases}
\begin{enumerate}[label=\alph*)]
\item \textbf{Convex $f$.} If $f$ is convex, the FW gap upper‑bounds the primal suboptimality:
\[
f(x_{k})-f^{\star}\le\gap(x_{k}).
\]
Hence (4) directly gives the classic $\mathcal{O}(1/K)$ convergence of the primal gap for the diminishing stepsize.

\item \textbf{Strongly Convex $f$ and Polyhedral $\dom$.} When $f$ is $\mu$‑strongly convex and $\dom$ is a polytope with pyramidal width $\delta>0$, the curvature constant $C_{f}$ satisfies $C_{f}\le L D^{2}\delta^{-1}$. Under the fixed stepsize $\gamma_{k}=\gamma\in(0,1)$, inequality (10) becomes a linear recursion,
\[
f(x_{k+1})-f^{\star}\le (1-\gamma\mu) \bigl(f(x_{k})-f^{\star}\bigr)+\frac{L D^{2}}{2}\gamma^{2},
\]
which yields linear convergence to a neighborhood of radius $\mathcal{O}(\gamma)$. Choosing $\gamma$ diminishing yields exact convergence.

\end{enumerate}

\end{document}